\documentclass[12pt, a4paper]{article}

% 1. Cấu hình Tiếng Việt và Font Times New Roman 13pt chuẩn
\usepackage[utf8]{inputenc}
\usepackage[T5]{fontenc}
\usepackage[vietnamese]{babel}
\usepackage{newtxtext,newtxmath}

\makeatletter
\renewcommand{\normalsize}{\@setfontsize\normalsize{13pt}{16pt}}
\makeatother
\normalsize

% 2. Gói Toán học & Ký hiệu bổ trợ
\usepackage{amsmath, amssymb, amsfonts, mathrsfs, amsthm}
\usepackage{graphicx}
\usepackage{fontawesome}
\usepackage{booktabs, tabularx, array, multirow}
\usepackage{enumitem}

% 3. Đồ họa TikZ, Bảng biến thiên & Biểu đồ
\usepackage{tikz}
\usetikzlibrary{calc, intersections, angles, quotes, patterns, arrows.meta, 3d}
\usepackage{pgfplots}
\pgfplotsset{compat=1.18}
\usepackage{tkz-euclide}
\usepackage{tkz-tab}

\renewcommand{\vec}[1]{\overrightarrow{#1}}

% 4. Hộp sư phạm (Tcolorbox)
\usepackage[many]{tcolorbox}
\tcbuselibrary{skins, breakable}

\newtcolorbox{pedabox}[1][]{
    enhanced,
    colback=blue!5!white,
    colframe=blue!75!black,
    fonttitle=\bfseries\sffamily,
    title=#1,
    boxrule=0.8pt,
    arc=2mm,
    breakable,
    left=3mm, right=3mm, top=2mm, bottom=2mm
}

\newtcolorbox{scaffoldbox}[1][]{
    enhanced,
    colback=orange!5!white,
    colframe=orange!85!black,
    fonttitle=\bfseries\sffamily,
    title=#1,
    boxrule=0.8pt,
    arc=2mm,
    breakable,
    left=3mm, right=3mm, top=2mm, bottom=2mm
}

\newtcolorbox{aibox}[1][]{
    enhanced,
    colback=green!5!white,
    colframe=teal!80!black,
    fonttitle=\bfseries\sffamily,
    title=#1,
    boxrule=0.8pt,
    arc=2mm,
    breakable,
    left=3mm, right=3mm, top=2mm, bottom=2mm
}

% 5. Căn lề chuẩn văn bản giáo án
\usepackage{geometry}
\geometry{
    a4paper,
    left=25mm,
    right=15mm,
    top=20mm,
    bottom=20mm
}

% 6. Header / Footer chuẩn hóa (BẮT BUỘC CHÍNH XÁC NỘI DUNG NÀY)
\usepackage{fancyhdr}
\usepackage{lastpage}
\pagestyle{fancy}
\fancyhf{}

\fancyhead[L]{\small \textit{Kế hoạch bài dạy Toán 11}}
\fancyhead[R]{\textbf{Trang \thepage/\pageref{LastPage}}}
\fancyfoot[L]{\small \textit{Tổ Toán - Tin}}
\fancyfoot[R]{\small \textit{Trường THCS\&THPT Hồng Vân}}
\renewcommand{\headrulewidth}{0.4pt}
\renewcommand{\footrulewidth}{0.4pt}

% 7. Giãn dòng & Giãn đoạn theo chuẩn Nghị định 30/2020/NĐ-CP
\usepackage{setspace}
\setstretch{1.15}
\setlength{\parindent}{1.0cm}
\setlength{\parskip}{5pt}

\begin{document}

\begin{center}
    {\Large \textbf{KẾ HOẠCH BÀI DẠY MÔN TOÁN LỚP 11}}\\[0.4em]
    {\LARGE \textbf{ÔN TẬP KIỂM TRA GIỮA HỌC KÌ I}}\\[0.5em]
    \textbf{Môn học:} Toán 11 | \textbf{Thời lượng:} 01 tiết (Tiết 22 theo PPCT)
\end{center}

\vspace{0.5em}
\hrule
\vspace{1em}

\section*{I. MỤC TIÊU DẠY HỌC}

\subsection*{1. Kiến thức, Năng lực}

\textbf{a) Năng lực Toán học đặc thù (nguyên văn chuẩn CT GDPT 2018):}
\begin{itemize}[leftmargin=1.5em]
    \item \textbf{Năng lực tư duy và lập luận toán học:}
    \begin{itemize}
        \item Hệ thống hóa và củng cố toàn diện các mạch kiến thức trọng tâm từ đầu năm học đến hết Chương III:
        \begin{itemize}
            \item \textbf{Chương I (Hàm số lượng giác và phương trình lượng giác):} Giá trị lượng giác của góc lượng giác; các hệ thức lượng giác cơ bản và công thức góc liên quan đặc biệt; công thức cộng, công thức nhân đôi, công thức biến đổi tích thành tổng và tổng thành tích; tập xác định, tính tuần hoàn, đồ thị của các hàm số lượng giác cơ bản; công thức nghiệm của các phương trình lượng giác cơ bản ($\sin x = m, \cos x = m, \tan x = m, \cot x = m$).
            \item \textbf{Chương II (Dãy số. Cấp số cộng và cấp số nhân):} Khái niệm dãy số, dãy số tăng, dãy số giảm, dãy số bị chặn; định nghĩa, số hạng tổng quát, tính chất và công thức tính tổng $n$ số hạng đầu tiên của cấp số cộng ($u_n = u_1 + (n - 1)d$, $S_n = \dfrac{n(u_1 + u_n)}{2}$) và cấp số nhân ($u_n = u_1 \cdot q^{n-1}$, $S_n = \dfrac{u_1(1 - q^n)}{1 - q}$).
            \item \textbf{Chương III (Các số đặc trưng đo xu thế trung tâm của mẫu số liệu ghép nhóm):} Mẫu số liệu ghép nhóm; cách tính số trung bình cộng ($\overline{x}$), trung vị ($M_e$), các tứ phân vị ($Q_1, Q_2, Q_3$), khoảng tứ phân vị ($\Delta_Q = Q_3 - Q_1$), mốt ($M_o$) và ý nghĩa thực tiễn của từng số đo.
        \end{itemize}
        \item So sánh, đối chiếu và nhận diện mối quan hệ biện chứng giữa các công thức đại số, giải tích và thống kê; phân tích và nhận diện nhanh các dạng bài tập trắc nghiệm và tự luận.
    \end{itemize}
    \item \textbf{Năng lực giải quyết vấn đề toán học:}
    \begin{itemize}
        \item Vận dụng linh hoạt các công thức để rút gọn biểu thức lượng giác, giải phương trình lượng giác có điều kiện nghiệm hoặc tìm số nghiệm thuộc khoảng, đoạn cho trước.
        \item Thiết lập và giải thành thạo hệ phương trình tìm các yếu tố cơ bản của cấp số cộng và cấp số nhân; tính tổng các dãy số quy về cấp số.
        \item Tính toán chính xác các số đặc trưng đo xu thế trung tâm của mẫu số liệu ghép nhóm; đọc hiểu và nhận xét bảng số liệu thực tế.
        \item Rèn luyện kĩ năng phân bổ thời gian hợp lí khi làm bài kiểm tra trắc nghiệm kết hợp tự luận (tỉ lệ 70\% trắc nghiệm khách quan, 30\% tự luận).
    \end{itemize}
    \item \textbf{Năng lực mô hình hóa toán học:}
    \begin{itemize}
        \item Giải quyết các bài toán thực tế tích hợp: mô hình dao động điều hòa trong Vật lí bằng hàm lượng giác, bài toán tăng trưởng dân số và lãi kép ngân hàng bằng cấp số, bài toán phân tích thể lực học sinh bằng mẫu số liệu ghép nhóm.
    \end{itemize}
    \item \textbf{Năng lực giao tiếp toán học:}
    \begin{itemize}
        \item Diễn đạt lời giải tự luận chặt chẽ, mạch lạc, đúng quy chuẩn toán học; sử dụng chính xác các kí hiệu toán học $u_n, d, q, S_n, \overline{x}, M_e, M_o, Q_k$.
    \end{itemize}
    \item \textbf{Năng lực sử dụng công cụ, phương tiện học toán:}
    \begin{itemize}
        \item Sử dụng thành thạo máy tính cầm tay (MTCT Casio fx-580VN X) để kiểm tra nghiệm phương trình lượng giác, giải hệ phương trình hai ẩn, tính tổng dãy số và xử lý các số đặc trưng thống kê một biến (\texttt{MENU} 6-1).
    \end{itemize}
\end{itemize}

\textbf{b) Năng lực số (theo Thông tư 02/2025/TT-BGDĐT):}
\begin{itemize}[leftmargin=1.5em]
    \item \textbf{Mã 2.6NC1a (Quản lý danh tính số và tài khoản an toàn):} Học sinh thực hành quản lý tài khoản ôn tập trực tuyến an toàn trên các nền tảng học tập số của nhà trường (Google Classroom, Azota, OLM); chủ động thiết lập mật khẩu mạnh, không tiết lộ mã xác thực hay chia sẻ tài khoản làm bài thi thử cho người khác, tuân thủ nghiêm ngặt quy chế kiểm tra đánh giá trên môi trường số.
\end{itemize}

\textbf{c) Năng lực AI (theo Quyết định 2422/QĐ-BGDĐT):}
\begin{itemize}[leftmargin=1.5em]
    \item \textbf{Mã 11.B2.1 (Đạo đức học tập và tính trung thực trong kỉ nguyên AI):} Học sinh nhận thức sâu sắc và cam kết thực hiện tính trung thực trong học tập; tuyệt đối không sử dụng các công cụ trí tuệ nhân tạo (AI) như ChatGPT, Gemini, PhotoMath để giải hộ đề cương kiểm tra hoặc làm bài gian lận; hiểu rằng AI chỉ đóng vai trò là gia sư gợi mở tư duy, không được phép thay thế năng lực tư duy toán học độc lập của bản thân.
\end{itemize}

\textbf{d) Năng lực chung \& Phẩm chất:}
\begin{itemize}[leftmargin=1.5em]
    \item \textbf{Năng lực chung:} Năng lực tự chủ và tự học (tự hệ thống kiến thức, chủ động rà soát các phần kiến thức còn yếu); giao tiếp và hợp tác nhóm khi cùng giải đề cương ôn tập.
    \item \textbf{Phẩm chất:} Trung thực, thẳng thắn trong học tập và kiểm tra; kiên trì, cẩn thận, có trách nhiệm với kết quả rèn luyện của bản thân.
\end{itemize}

\section*{II. THIẾT BỊ DẠY HỌC VÀ HỌC LIỆU}
\begin{itemize}[leftmargin=1.5em]
    \item \textbf{Giáo viên:} Kế hoạch bài dạy chuẩn CV 5512; bài trình chiếu PowerPoint hệ thống hóa ma trận kiến thức giữa kì I; Phiếu học tập số 1 (Ma trận công thức 3 chương), Phiếu học tập số 2 (Đề cương bài tập ôn tập chuẩn hóa phân tầng); máy tính, máy chiếu.
    \item \textbf{Học sinh:} Vở ghi bài, SGK Toán 11; máy tính cầm tay Casio fx-580VN X; đề cương ôn tập giữa kì I đã chuẩn bị trước ở nhà.
\end{itemize}

\section*{III. TIẾN TRÌNH DẠY HỌC (01 TIẾT --- 45 PHÚT)}

\begin{center}
    \framebox[1.0\textwidth]{\parbox{0.96\textwidth}{
        \textbf{CẤU TRÚC TIẾN TRÌNH TIẾT 22 (ÔN TẬP KIỂM TRA GIỮA HỌC KÌ I):}\\[0.3em]
        $\bullet$ \textbf{Hoạt động 1 (05 phút):} Khởi động --- Trò chơi "Vượt chướng ngại vật" quét nhanh các công thức chìa khóa 3 chương.\\[0.3em]
        $\bullet$ \textbf{Hoạt động 2 (12 phút):} Hệ thống hóa kiến thức trọng tâm \& Cấu trúc ma trận đề thi --- Giàn giáo sư phạm HS TB - Yếu --- Tích hợp Năng lực số (Mã 2.6NC1a) \& Năng lực AI (Mã 11.B2.1).\\[0.3em]
        $\bullet$ \textbf{Hoạt động 3 (20 phút):} Luyện tập giải đề minh họa phân hóa (3 bài toán đại diện cho 3 chương).\\[0.3em]
        $\bullet$ \textbf{Hoạt động 4 (08 phút):} Vận dụng & Hướng dẫn chiến thuật làm bài thi trắc nghiệm kết hợp tự luận.
    }}
\end{center}

\vspace{0.8em}
\hrule
\vspace{0.8em}

\subsubsection*{1. Hoạt động 1: Mở đầu / Khởi động (05 phút)}
\textbf{Chủ đề:} \textit{Trò chơi "Vượt chướng ngại vật": 4 câu hỏi chớp nhoáng}
\begin{itemize}
    \item \textbf{a) Mục tiêu:} Kích hoạt trí nhớ dài hạn của học sinh về các công thức then chốt của 3 chương đã học; tạo tâm thế sẵn sàng, tự tin bước vào tiết ôn tập.
    \item \textbf{b) Nội dung:} Giáo viên chiếu 4 câu hỏi trắc nghiệm nhanh lên màn hình:
    \begin{itemize}
        \item \textbf{Câu 1 (Chương I):} Phương trình $\sin x = \sin \alpha$ có các nghiệm là gì?
        \item \textbf{Câu 2 (Chương I):} Đẳng thức $\sin 2x = 2\sin x \cos x$ đúng hay sai?
        \item \textbf{Câu 3 (Chương II):} Cấp số cộng có $u_1 = 2, d = 3$ thì số hạng thứ $5$ ($u_5$) bằng bao nhiêu?
        \item \textbf{Câu 4 (Chương III):} Để tìm trung vị của mẫu số liệu ghép nhóm có $n = 40$, ta tìm nhóm đầu tiên có tần số tích lũy đạt mốc nào?
    \end{itemize}
    \item \textbf{c) Sản phẩm:} Câu trả lời của học sinh:
    \begin{itemize}
        \item Câu 1: $x = \alpha + k2\pi$ hoặc $x = \pi - \alpha + k2\pi$ ($k \in \mathbb{Z}$).
        \item Câu 2: Đúng (công thức góc nhân đôi).
        \item Câu 3: $u_5 = u_1 + 4d = 2 + 4 \cdot 3 = 14$.
        \item Câu 4: Đạt mốc $\dfrac{n}{2} = \dfrac{40}{2} = 20$.
    \end{itemize}
    \item \textbf{d) Tổ chức thực hiện:} GV lần lượt chiếu câu hỏi $\to$ HS giơ thẻ đáp án A, B, C, D hoặc trả lời đồng thanh $\to$ GV nhận xét, ghi nhận điểm cộng và dẫn dắt vào bài ôn tập.
\end{itemize}

\subsubsection*{2. Hoạt động 2: Hệ thống hóa kiến thức trọng tâm \& Cấu trúc ma trận đề thi (12 phút)}

\paragraph{Mục 1: Bảng tổng hợp các trục kiến thức then chốt giữa học kì I}

\begin{center}
\begin{tabularx}{\linewidth}{|l|X|X|}
\hline
\textbf{Chương học} & \textbf{Kiến thức cốt lõi} & \textbf{Kĩ năng toán học cần đạt} \\
\hline
\textbf{Chương I} & $\bullet$ Công thức lượng giác (cộng, nhân đôi, biến đổi).\\
(Lượng giác) & $\bullet$ Tập xác định, chu kì tuần hoàn các hàm số.\\
& $\bullet$ Phương trình cơ bản: $\sin, \cos, \tan, \cot$. & $\bullet$ Rút gọn biểu thức lượng giác.\\
& & $\bullet$ Giải phương trình lượng giác cơ bản và tìm nghiệm thỏa mãn điều kiện. \\
\hline
\textbf{Chương II} & $\bullet$ Khái niệm dãy số (tăng, giảm, bị chặn).\\
(Cấp số) & $\bullet$ Cấp số cộng: $u_n = u_1 + (n - 1)d$; $S_n$.\\
& $\bullet$ Cấp số nhân: $u_n = u_1 \cdot q^{n-1}$; $S_n$. & $\bullet$ Giải hệ phương trình tìm $u_1, d$ và $u_1, q$.\\
& & $\bullet$ Nhận biết dãy số, tính tổng $n$ số hạng đầu; giải bài toán thực tế lãi kép, tăng trưởng. \\
\hline
\textbf{Chương III} & $\bullet$ Bảng tần số, tần số tương đối ghép nhóm.\\
(Thống kê) & $\bullet$ Số trung bình $\overline{x}$; Mốt $M_o$.\\
& $\bullet$ Trung vị $M_e$; Tứ phân vị $Q_1, Q_3$; Khoảng $\Delta_Q$. & $\bullet$ Lập bảng 4 cột tính số trung bình $\overline{x}$.\\
& & $\bullet$ Xác định nhóm chứa phân vị qua tần số tích lũy $cf_i$ và tính chính xác $M_e, M_o, Q_k$. \\
\hline
\end{tabularx}
\end{center}

\begin{scaffoldbox}[Giàn giáo sư phạm cho HS TB - Yếu: Mẹo tránh "bẫy" đề thi và phân bổ thời gian]
    \textbf{1. Phân bổ thời gian làm bài thi 90 phút (Trắc nghiệm 70\% + Tự luận 30\%):}\\
    $\bullet$ \textbf{Phần Trắc nghiệm (35 câu - 7,0 điểm):} Dành tối đa $45 - 50$ phút. Làm từ câu dễ đến câu khó; câu nào chưa ra đáp án sau $2$ phút thì đánh dấu tạm thời và chuyển sang câu kế tiếp.\\
    $\bullet$ \textbf{Phần Tự luận (3 - 4 bài - 3,0 điểm):} Dành $35$ phút. Trình bày cẩn thận từng bước giải, không bỏ sót điều kiện xác định (ví dụ điều kiện của hàm $\tan, \cot$ hay mẫu số khác $0$). Dành $5$ phút cuối giờ soát lại mã đề và số báo danh.\\
    \textbf{2. Ba lỗi sai kinh điển cần tránh:}\\
    $\bullet$ Nghiệm phương trình lượng giác: Quên đuôi chu kì $+ k2\pi$ hoặc $+ k\pi$; quên điều kiện $k \in \mathbb{Z}$.\\
    $\bullet$ Cấp số: Nhầm giữa công sai $d$ (cộng) và công bội $q$ (nhân); nhầm số mũ $q^{n-1}$ thành $q^n$.\\
    $\bullet$ Thống kê: Nhầm nhóm chứa trung vị do nhìn nhầm cột tần số $m_i$ thay vì cột tần số tích lũy $cf_i$.
\end{scaffoldbox}

\paragraph{Mục 2: Tích hợp Năng lực số (Mã 2.6NC1a) \& Năng lực AI (Mã 11.B2.1)}

\begin{pedabox}[Năng lực số (Mã 2.6NC1a): Bảo mật tài khoản ôn tập trực tuyến]
    $\bullet$ Giáo viên hướng dẫn học sinh đăng nhập vào hệ thống ngân hàng đề thi trực tuyến của trường; nhắc nhở học sinh không lưu mật khẩu trên máy tính công cộng tại quán net hay máy tính lạ.\\
    $\bullet$ Học sinh đổi mật khẩu định kì và bật xác thực bảo mật tài khoản học tập cá nhân.
\end{pedabox}

\begin{aibox}[Tích hợp Năng lực AI (QĐ 2422/QĐ-BGDĐT --- Mã 11.B2.1): Trung thực trong học tập]
    $\bullet$ \textbf{Cam kết văn hóa học tập trung thực:}\\
    Trước kì kiểm tra giữa kì I, giáo viên quán triệt tinh thần tự lực: Các ứng dụng AI giải toán chụp ảnh có thể đưa ra đáp số nhanh nhưng không thể thay thế kiến thức trong đầu học sinh khi bước vào phòng thi trực tiếp trên giấy.\\
    $\bullet$ Học sinh cam kết không dùng AI để gian lận khi làm bài kiểm tra thử, dùng năng lực thực chất để rà soát chính xác lỗ hổng kiến thức của bản thân.
\end{aibox}

\subsubsection*{3. Hoạt động 3: Luyện tập giải đề minh họa phân hóa (20 phút)}

\paragraph{Dạng 1: Bài toán Lượng giác trọng tâm (06 phút)}
\begin{itemize}
    \item \textbf{Bài toán 1:} Giải phương trình lượng giác sau:
    $$2\cos\left(2x + \dfrac{\pi}{3}\right) + \sqrt{3} = 0$$
    Tìm tất cả các nghiệm của phương trình thuộc khoảng $(0; \pi)$.
    \item \textit{Lời giải chi tiết:}\\
    Phương trình tương đương:
    $$\cos\left(2x + \dfrac{\pi}{3}\right) = -\dfrac{\sqrt{3}}{2} \iff \cos\left(2x + \dfrac{\pi}{3}\right) = \cos\left(\dfrac{5\pi}{6}\right)$$
    $$\iff \begin{bmatrix} 2x + \dfrac{\pi}{3} = \dfrac{5\pi}{6} + k2\pi \\ 2x + \dfrac{\pi}{3} = -\dfrac{5\pi}{6} + k2\pi \end{bmatrix} \iff \begin{bmatrix} 2x = \dfrac{\pi}{2} + k2\pi \\ 2x = -\dfrac{7\pi}{6} + k2\pi \end{bmatrix} \iff \begin{bmatrix} x = \dfrac{\pi}{4} + k\pi \\ x = -\dfrac{7\pi}{12} + k\pi \end{bmatrix} \quad (k \in \mathbb{Z})$$
    Xét các nghiệm thuộc khoảng $(0; \pi)$:\\
    - Với họ nghiệm $x = \dfrac{\pi}{4} + k\pi$:
    $$0 < \dfrac{\pi}{4} + k\pi < \pi \iff -\dfrac{1}{4} < k < \dfrac{3}{4} \implies k = 0 \implies x_1 = \dfrac{\pi}{4}$$
    - Với họ nghiệm $x = -\dfrac{7\pi}{12} + k\pi$:
    $$0 < -\dfrac{7\pi}{12} + k\pi < \pi \iff \dfrac{7}{12} < k < \dfrac{19}{12} \implies k = 1 \implies x_2 = -\dfrac{7\pi}{12} + \pi = \dfrac{5\pi}{12}$$
    Vậy phương trình có hai nghiệm thuộc khoảng $(0; \pi)$ là $x = \dfrac{\pi}{4}$ và $x = \dfrac{5\pi}{12}$.
\end{itemize}

\paragraph{Dạng 2: Bài toán Cấp số cộng \& Cấp số nhân (07 phút)}
\begin{itemize}
    \item \textbf{Bài toán 2:} Cho cấp số cộng $(u_n)$ thỏa mãn hệ phương trình:
    $$\begin{cases} u_1 + u_5 = 14 \\ u_2 + u_6 = 18 \end{cases}$$
    \begin{itemize}
        \item a) Tìm số hạng đầu $u_1$ và công sai $d$ của cấp số cộng.
        \item b) Tính tổng của $25$ số hạng đầu tiên ($S_{25}$).
    \end{itemize}
    \item \textit{Lời giải chi tiết:}\\
    a) Biểu diễn các số hạng qua $u_1$ và $d$ ($u_k = u_1 + (k - 1)d$):
    $$\begin{cases} u_1 + (u_1 + 4d) = 14 \\ (u_1 + d) + (u_1 + 5d) = 18 \end{cases} \iff \begin{cases} 2u_1 + 4d = 14 \\ 2u_1 + 6d = 18 \end{cases}$$
    Trừ phương trình dưới cho phương trình trên:
    $$2d = 4 \iff d = 2$$
    Thay $d = 2$ vào phương trình thứ nhất:
    $$2u_1 + 4(2) = 14 \iff 2u_1 = 6 \iff u_1 = 3$$
    Vậy cấp số cộng có số hạng đầu $u_1 = 3$ và công sai $d = 2$.\\
    b) Áp dụng công thức tính tổng $n$ số hạng đầu tiên:
    $$S_{25} = \dfrac{25 \cdot [2u_1 + (25 - 1)d]}{2} = \dfrac{25 \cdot [2 \cdot 3 + 24 \cdot 2]}{2} = \dfrac{25 \cdot (6 + 48)}{2} = \dfrac{25 \cdot 54}{2} = 25 \cdot 27 = 675$$
\end{itemize}

\paragraph{Dạng 3: Bài toán Thống kê mẫu số liệu ghép nhóm (07 phút)}
\begin{itemize}
    \item \textbf{Bài toán 3:} Khảo sát điểm kiểm tra đánh giá thường xuyên môn Toán của một lớp gồm $45$ học sinh được ghi nhận ở bảng ghép nhóm sau:
    \begin{center}
    \begin{tabular}{|c|c|c|c|c|}
    \hline
    \textbf{Nhóm điểm} & $[4; 6)$ & $[6; 7)$ & $[7; 8)$ & $[8; 10]$ \\
    \hline
    \textbf{Số học sinh ($m_i$)} & $9$ & $15$ & $13$ & $8$ \\
    \hline
    \end{tabular}
    \end{center}
    \begin{itemize}
        \item a) Tính số trung bình cộng $\overline{x}$ của mẫu số liệu ghép nhóm.
        \item b) Xác định nhóm chứa trung vị và tính trung vị $M_e$.
    \end{itemize}
    \item \textit{Lời giải chi tiết:}\\
    a) Xác định giá trị đại diện của từng nhóm:
    $c_1 = \dfrac{4 + 6}{2} = 5$; $c_2 = \dfrac{6 + 7}{2} = 6{,}5$; $c_3 = \dfrac{7 + 8}{2} = 7{,}5$; $c_4 = \dfrac{8 + 10}{2} = 9$.\\
    Số trung bình cộng của lớp là:
    $$\overline{x} = \dfrac{9 \cdot 5 + 15 \cdot 6{,}5 + 13 \cdot 7{,}5 + 8 \cdot 9}{45} = \dfrac{45 + 97{,}5 + 97{,}5 + 72}{45} = \dfrac{312}{45} \approx 6{,}93$$
    b) Bảng tần số tích lũy: $cf_1 = 9$; $cf_2 = 9 + 15 = 24$; $cf_3 = 24 + 13 = 37$; $cf_4 = 45$.\\
    Ta có: $\dfrac{n}{2} = \dfrac{45}{2} = 22{,}5$.\\
    Nhóm đầu tiên có tần số tích lũy $cf \ge 22{,}5$ là nhóm $[6; 7)$ (vì $cf_2 = 24 \ge 22{,}5$).\\
    Các thông số: $a_2 = 6$, độ dài nhóm $h = 7 - 6 = 1$, tần số $m_2 = 15$, tần số tích lũy nhóm trước $cf_1 = 9$.\\
    Áp dụng công thức tính trung vị:
    $$M_e = 6 + \dfrac{22{,}5 - 9}{15} \cdot 1 = 6 + \dfrac{13{,}5}{15} = 6 + 0{,}9 = 6{,}9$$
\end{itemize}

\subsubsection*{4. Hoạt động 4: Vận dụng \& Hướng dẫn làm bài thi (08 phút)}
\textbf{Chủ đề:} \textit{Chiến thuật làm bài thi đạt kết quả tối ưu}
\begin{itemize}
    \item \textbf{a) Mục tiêu:} Hệ thống hóa phương pháp làm bài trắc nghiệm kết hợp tự luận; hướng dẫn sử dụng MTCT hỗ trợ kiểm tra kết quả nhanh chóng.
    \item \textbf{b) Nội dung:}
    \begin{itemize}
        \item Kĩ thuật kiểm tra nghiệm phương trình lượng giác bằng lệnh \texttt{CALC} trên MTCT Casio fx-580VN X: Chuyển máy về chế độ Radian (\texttt{SHIFT} $\to$ \texttt{MENU} $\to$ \texttt{2} $\to$ \texttt{2}), nhập vế trái trừ vế phải, bấm \texttt{CALC} nhập từng giá trị của các phương án trắc nghiệm.
        \item Kĩ thuật kiểm tra tổng cấp số bằng lệnh $\sum$ (\texttt{SHIFT} $x$).
    \end{itemize}
    \item \textbf{c) Sản phẩm:} Học sinh nắm vững thao tác máy tính cầm tay để tự tin kiểm tra lại bài làm trong phòng thi; hoàn thiện mục tiêu cá nhân cho kì kiểm tra sắp tới.
    \item \textbf{d) Tổ chức thực hiện:} GV thao tác mẫu trên máy chiếu $\to$ HS thực hành đồng loạt tại chỗ $\to$ GV dặn dò ôn tập theo đề cương và phổ biến quy chế thi.
\end{itemize}

\newpage

\section*{IV. PHỤ LỤC HỌC LIỆU BỔ TRỢ}

\subsection*{Phụ lục 1: Ma trận đề kiểm tra giữa học kì I môn Toán 11}

\begin{center}
\begin{tabularx}{\linewidth}{|p{3.5cm}|c|c|c|c|c|}
\hline
\textbf{Chủ đề kiến thức} & \textbf{Nhận biết} & \textbf{Thông hiểu} & \textbf{Vận dụng} & \textbf{Vận dụng cao} & \textbf{Tổng số câu} \\
\hline
\textbf{Chương I: Lượng giác} & 6 câu TN & 4 câu TN & 2 câu TN + 1 TL & 1 câu TL & 12 TN + 2 TL \\
\hline
\textbf{Chương II: Dãy số, Cấp số} & 6 câu TN & 4 câu TN & 2 câu TN + 1 TL & 1 câu TL & 12 TN + 2 TL \\
\hline
\textbf{Chương III: Thống kê ghép nhóm} & 4 câu TN & 4 câu TN & 3 câu TN + 1 TL & --- & 11 TN + 1 TL \\
\hline
\textbf{Tổng cộng} & \textbf{16 câu TN} & \textbf{12 câu TN} & \textbf{7 TN + 3 TL} & \textbf{2 câu TL} & \textbf{35 TN + 5 TL} \\
\hline
\textbf{Tỉ lệ điểm số} & \textbf{3,2 điểm} & \textbf{2,4 điểm} & \textbf{3,0 điểm} & \textbf{1,4 điểm} & \textbf{10,0 điểm} \\
\hline
\end{tabularx}
\end{center}

\subsection*{Phụ lục 2: Phiếu học tập --- Đề cương ôn tập tự luyện chắt lọc}
\noindent\textbf{Họ và tên:} .................................................................... \textbf{Lớp:} 11A....... \textbf{Nhóm:} ..........

\begin{pedabox}[Một số câu hỏi rèn luyện then chốt]
    $\bullet$ \textbf{Câu 1 (Lượng giác):} Cho $\cos\alpha = -\dfrac{3}{5}$ với $\dfrac{\pi}{2} < \alpha < \pi$. Tính $\sin 2\alpha$ và $\cos 2\alpha$.\\
    $\bullet$ \textbf{Câu 2 (Cấp số cộng):} Cho cấp số cộng có $u_4 = 10, u_7 = 19$. Tìm số hạng đầu $u_1$ và tổng $S_{20}$.\\
    $\bullet$ \textbf{Câu 3 (Cấp số nhân):} Tìm $x$ để ba số $x - 2, x + 1, x + 7$ theo thứ tự lập thành một cấp số nhân.\\
    $\bullet$ \textbf{Câu 4 (Thống kê):} Cho mẫu số liệu ghép nhóm có $n = 60$, nhóm $[50; 60)$ có tần số $m_3 = 20$, tần số tích lũy nhóm trước $cf_2 = 18$. Tính trung vị $M_e$.
\end{pedabox}

\subsection*{Phụ lục 3: Bảng Rubric đánh giá năng lực trong Tiết ôn tập giữa kì I}

\begin{center}
\begin{tabularx}{\linewidth}{|p{3.2cm}|X|X|X|X|}
\hline
\textbf{Tiêu chí} & \textbf{Mức 1 (Chưa đạt)} & \textbf{Mức 2 (Đạt)} & \textbf{Mức 3 (Khá)} & \textbf{Mức 4 (Tốt)} \\
\hline
\textbf{1. Hệ thống hóa kiến thức 3 chương} & Quên nhiều công thức cốt lõi; chưa giải được các bài toán cơ bản. & Nhớ các công thức cơ bản; giải đúng phương trình lượng giác và cấp số đơn giản. & Nắm chắc toàn bộ kiến thức 3 chương; giải thành thạo bài toán vận dụng. & Tổng hợp kiến thức xuất sắc; biến đổi linh hoạt và giải quyết tốt các câu hỏi phân hóa cao. \\
\hline
\textbf{2. Kĩ năng số \& MTCT (Mã 2.6NC1a)} & Chưa biết cách bấm máy kiểm tra nghiệm hoặc quên mật khẩu tài khoản học. & Đăng nhập tài khoản học an toàn; dùng MTCT giải hệ phương trình. & Thành thạo các tính năng bảng tính và kiểm tra nghiệm lượng giác trên MTCT. & Quản lý danh tính số chuẩn mực; tối ưu hóa công cụ số để kiểm tra bài thi chính xác. \\
\hline
\textbf{3. Đạo đức học tập AI (Mã 11.B2.1)} & Còn phụ thuộc vào AI giải bài hộ; thiếu tự giác ôn tập. & Nhận thức được yêu cầu trung thực trong thi cử và kiểm tra. & Nghiêm túc tự làm đề cương, cam kết không gian lận bằng AI trong phòng thi. & Lan tỏa tinh thần học tập trung thực, tự chủ và làm chủ công nghệ đến toàn thể lớp học. \\
\hline
\end{tabularx}
\end{center}

\end{document}
